\bookchapter{Bab 11: Fotografi dan manajemen gambar}

\chapterintro{1. Apakah Anda Juga Pernah Membuat Kesalahan Ini?}

\booksection{Pertama mari bicara tentang beberapa masalah paling umum, lihat berapa banyak yang Anda alami:}

Lensa Kotor

\begin{itemize}
  \item \textbf{: Ini benar-benar masalah paling umum. Saat mengeluarkan ponsel, sidik jari menempel di lensa, hasil foto berkabut, dikira ponsel rusak. Biasakan, sebelum memotret lap lensa dengan ujung baju.}Tidak Pernah Mengatur Eksposur
  \item \textbf{: Dalam mode default, ketuk layar adalah fokus, tetapi jika setelah mengetuk Anda geser ke atas dan bawah, Anda bisa mengatur kecerahan. Banyak orang memotret terlalu gelap atau terlalu terang, karena tidak tahu operasi ini.}Lampu Kilat Selalu Menyala
  \item \textbf{: Hasil foto lampu kilat LED ponsel benar-benar jelek, cahaya keras dan dingin, membuat orang terlihat seperti hantu. Kecuali benar-benar gelap gulita, matikan saja. Dalam kondisi cahaya redup, mode malam jauh lebih baik daripada lampu kilat.}Zoom dengan Mengepit Dua Jari
  \item \textbf{: Pinch-to-zoom di layar adalah zoom digital, pada dasarnya memotong gambar lalu memperbesar, kualitas gambar jelas menurun. Jika ingin memotret objek jauh, langsung gunakan tombol zoom di viewfinder (0.5x/1x/2x/3x), ini adalah zoom optik, tidak mengurangi kualitas gambar.}: Pinch zoom di layar adalah zoom digital, pada dasarnya memotong gambar lalu memperbesar, kualitas gambar akan menurun jelas. Jika ingin memotret jarak jauh, langsung gunakan tombol zoom di viewfinder (0.5x/1x/2x/3x). Di antaranya 0.5x (ultra wide) dan 1x (kamera utama) adalah lensa optik native, tidak mengurangi kualitas; 2x/3x kebanyakan adalah "zoom kualitas optik" — dicapai melalui pemotongan sensor piksel tinggi, kualitas mendekati zoom optik tapi bukan lensa independen. Zoom di atas 3x biasanya adalah zoom hybrid atau zoom digital, kualitas gambar akan menurun nyata.
\end{itemize}

\booksection{Ini tidak memerlukan pengetahuan teori, gunakan sekali langsung melihat hasilnya:}

Rule of Thirds + Grid Lines

\begin{itemize}
  \item \textbf{: Aktifkan "Grid Lines" (juga disebut reference lines) di pengaturan kamera. Saat memotret, letakkan objek yang ingin Anda tonjolkan di perpotongan empat garis, komposisi langsung terasa "pas". Ini adalah aturan paling sederhana dan paling efektif dalam fotografi, bukan ilmu gaib, hanya sesuai dengan kebiasaan visual mata manusia.}Cari "Leading Lines"
  \item \textbf{: Saat memotret pemandangan atau street view, perhatikan apakah ada garis alami dalam frame—jalan, pagar, dinding, deretan pohon—biarkan garis-garis ini membentang dari tepi frame ke subjek. Jika trik ini dikuasai, orang yang melihat foto Anda akan mengira Anda pernah belajar fotografi.}Jangan Selalu Berdiri Saat Memotret
  \item \textbf{: Kebanyakan orang memotret selalu dari ketinggian mata sendiri, jadi sudut pandang hasilnya monoton. Coba jongkok memotret dari sudut rendah, atau cari tempat tinggi memotret ke bawah, suasana yang dihasilkan dari tempat yang sama benar-benar berbeda. Memotret anak dan hewan peliharaan terutama perlu jongkok, sejajar mata baru bisa menghasilkan rasa alami.}Cahaya Adalah Faktor Utama
  \item \textbf{: Memotret portrait di bawah matahari terik siang hari, bayangan di wajah keras dan jelek. Cahaya terbaik adalah "golden hour" pagi dan sore, matahari miring, warna cenderung hangat, memotret apa pun terlihat bagus. Cahaya mendung meski tidak menonjol, tetapi menang karena merata dan lembut, memotret portrait justru lebih mudah menghasilkan foto bagus daripada di bawah matahari terik.}Tips Penggunaan Mode Malam
  \item \textbf{: Setelah membuka mode malam, ponsel akan otomatis memotret beberapa frame lalu menggabungkannya. Saat ini yang paling penting—tangan tidak boleh gemetar. Bisa bersandar ke dinding, bisa diletakkan di meja, jika tidak bisa tahan napas. Mode malam memerlukan Anda stabil selama 2-4 detik, jika gemetar gambar akan blur.}3. Mode Portrait dan Mode Pro
\end{itemize}

\booksection{Mode Portrait}

\begin{itemize}
  \item \textbf{: Dapat menghasilkan efek background blur, subjek jelas (depth of field dangkal). Karena ukuran fisik lensa ponsel kecil, kemampuan bokeh asli terbatas, jadi mengandalkan algoritma AI—sederhananya AI memotong Anda keluar, lalu mengaburkan background. Saat digunakan pertahankan jarak 1-2 meter untuk hasil terbaik, terlalu jauh atau terlalu dekat algoritma mudah salah. Juga, saat memotret portrait, perhatikan background jangan ada garis terlalu kompleks, jika tidak AI masking akan "bermasalah" di tepi rambut.}Mode Pro
  \item \textbf{: Jika Anda ingin kendali penuh atas efek gambar, coba mode pro. Beberapa parameter kunci perlu diketahui: ISO makin rendah makin bersih kualitas gambar, tetapi gambar jadi gelap; shutter makin lambat makin banyak cahaya masuk, tetapi tangan gemetar akan blur; white balance mengatur suhu warna, jika tidak diatur dengan baik foto akan kebiruan atau kekuningan. Orang biasa tidak perlu terus-menerus menggunakan mode pro, tetapi memahami beberapa parameter ini setidaknya membuat Anda tahu apa yang dilakukan kamera.}4. Editing Foto: Jangan Berlebihan, Tepat Saja Sudah Cukup
\end{itemize}

\booksection{Banyak orang salah paham tentang editing, menganggap hasil edit jadi tidak asli. Sebenarnya foto yang diambil kamera memang sudah berbeda dengan yang dilihat mata telanjang, penyesuaian yang tepat adalah untuk mengembalikan apa yang Anda lihat saat itu, bukan pemalsuan.}

Snapseed

\begin{itemize}
  \item \textbf{: Alat editing gratis dari Google, fitur profesional tanpa iklan. Alur editing harian saya: pertama gunakan "Tune Image" untuk menarik kecerahan dan kontras gunakan "Details" untuk menajamkan gunakan "White Balance" untuk mengoreksi warna terakhir "Crop" ke komposisi yang diinginkan. Fitur "Selective" Snapseed sangat berguna, bisa hanya menyesuaikan area tertentu dalam gambar, tanpa mengubah seluruh gambar.}Xingtu/Meitu
  \item \textbf{: Keduanya cocok untuk editing portrait. Tapi menurut saya ada prinsip: filter beauty dibawah 30\% sudah cukup, lebih dari itu jadi tidak natural, kulit seperti plastik. Natural lebih bagus daripada sempurna.}Lightroom Mobile
  \item \textbf{: Jika Anda menuntut soal warna, ini yang paling profesional. Memiliki penyesuaian kurva, HSL per-warna, dan fitur lanjutan lainnya. Namun bagi kebanyakan orang agak berlebihan.}Saran Alur Editing
  \item \textbf{: Kebiasaan saya sendiri adalah memotret foto asli dengan baik, setelah pulang perlahan mengedit di layar komputer atau iPad, layar ponsel terlalu kecil, warna juga kurang akurat. Jika harus mengedit di ponsel, ingat atur kecerahan layar ke menengah-atas, jangan mengedit dalam mode hemat baterai.}5. Foto Terlalu Banyak? Strategi Pengorganisasian
\end{itemize}

\booksection{Album ponsel menumpuk ribuan foto, mencari satu foto tertentu butuh waktu lama—pengalaman ini semua orang punya. Beberapa metode pengorganisasian praktis:}

Manfaatkan Pencarian

\begin{itemize}
  \item \textbf{: Sekarang album iPhone dan kebanyakan ponsel Android mendukung pencarian cerdas. Anda bisa mencari "kucing" "pantai" "KTP" "bunga" kata kunci ini, sistem akan menggunakan AI untuk mengenali konten gambar dan membantu menemukannya. Juga bisa mencari tanggal dan lokasi, seperti "Mei tahun lalu" "Beijing". Fitur ini tidak banyak yang tahu, tetapi benar-benar sangat berguna.}Bersihkan Sampah Secara Berkala
  \item \textbf{: Screenshot, foto duplikat, foto gagal yang blur, inilah biang kerok yang menghabiskan ruang. Album iOS memiliki fitur "Duplicate Items" otomatis merge; Android bisa cari folder "screenshots" di file manager, hapus batch screenshot yang tidak diperlukan. Saya menghabiskan lima menit setiap bulan untuk membersihkan, bisa mengosongkan beberapa GB.}Buat Album, Jangan Hanya Andalkan Sortir Default
  \item \textbf{: Sortir default sistem berdasarkan waktu tidak salah, tetapi untuk konten lintas waktu seperti "foto tumbuh kembang anak" "foto perjalanan" "dokumen penting", Anda harus membuat album manual untuk mengkategorikan. "Shared Album" Apple juga bisa menambahkan foto bersama keluarga, satu orang membuat, seluruh keluarga bisa melihat.}Pencadangan Adalah Garis Bawah
  \item \textbf{: Foto adalah data yang paling tidak tergantikan—file hilang bisa dibuat ulang, riwayat chat hilang ya sudahlah, tetapi foto anak, video orang tua hilang benar-benar hilang. Wajib aktifkan pencadangan cloud. iCloud Photos paling nyaman, Android berbagai merek punya cloud service sendiri, atau gunakan pencadangan otomatis Baidu Netdisk/Aliyun Drive. Prinsipnya: setidaknya ada satu cadangan cloud, jangan taruh semua telur di satu keranjang.}6. Screenshot dan Screen Recording
\end{itemize}

\booksection{Screenshot}

\begin{itemize}
  \item \textbf{: Tekan bersamaan tombol power + volume up, metode universal semua ponsel. Banyak ponsel Android juga mendukung screenshot tiga jari swipe ke bawah, screenshot ketukan dua ruas jari. Long screenshot (scrolling screenshot) di Android sudah matang—setelah screenshot satu ketuk tombol "Scroll", ponsel akan otomatis scroll ke bawah dan menggabungkan menjadi satu gambar panjang, cocok untuk screenshot riwayat chat dan halaman web.}Screen Recording
  \item \textbf{: Tambahkan "Screen Recording" dari Control Center, ketuk sekali untuk mulai merekam, ketuk lagi untuk berhenti. Sebelum merekam tekan lama tombol rekam untuk memilih apakah merekam suara mikrofon—jika ingin membuat tutorial narasi nyalakan mic, jika hanya merekam operasi layar matikan. Video yang selesai direkam otomatis disimpan ke album.}Berlatih: Selesaikan Satu Sesi Foto + Edit + Organisir dengan Ponsel
\end{itemize}

\begin{tcbpractice}{Aktifkan "Grid Lines" di pengaturan kamera, cari tempat dengan cahaya alami ambil tiga foto (perhatikan lap lensa, atur eksposur).}
\begin{enumerate}
  \item Pilih yang terbaik, gunakan Snapseed untuk penyesuaian dasar (kecerahan, kontras, crop).
  \item Buat album khusus untuk foto ini, beri nama "Karya Saya".
  \item Periksa apakah pencadangan cloud telah diaktifkan, pastikan foto sudah tersimpan satu salinan di cloud.
  \item Lampiran Bab 11: Pengetahuan Teknik Imaging Profesional
\end{enumerate}

\end{tcbpractice}
